\documentclass[12pt]{article}
\usepackage{notes-common}
\geometry{margin=0.35in}
\AtBeginDocument{\fontsize{14}{17}\selectfont}

% Fri Sep 25 -- Electric potential energy. Same lesson for both tracks
% (no calculus needed), so no PHYS 230/235 split.
%
% Restructured 2026-09-24 per Seth: gravity (height, then a hole) first,
% then a uniform E field for BOTH + and - charges, then the full
% (non-uniform) Newtonian form of gravitational PE with its own graph
% ("higher is more PE" even though U<0 everywhere), then electric PE of
% point charges (with a blank set of axes right after U(q_1,q_2) for
% sketching U vs. r for both opposite and like charges), and finally the
% dipole. See electric-pe-notes-key.tex for the filled-in, all-diagrams,
% all-curves version -- every diagram and graph below is the SAME axes/
% frame as the key, just without the student-drawn content.
%
% Print layout matches the redesigned Gauss's-law handout: small margins,
% bigger base font, no ruled writing lines, no task/derivation boxes ---
% just blank space. Major standalone equations are blanked as a labeled
% line followed by its own underline (not an inline blank right after
% "="), since that inline style doesn't leave room for a fraction.
%
% \eqline[gap]{label} -- [gap] is the blank space between the label and the
% rule, i.e. the room a student actually has to write the answer. Default
% (-0.5mm) suits a short answer (mgd, qEd, -pE\cos\theta). A single-level
% fraction (U(r)=-GMm/r) or worse, a two-level one / a sum-with-limits
% (U(q_1,q_2), U_total) needs real room -- measured directly off
% electric-pe-notes-key.pdf at 300dpi: U(r) and U(q_1,q_2) both typeset at
% ~0.41-0.43in tall, U_total (fraction AND a sum) at ~0.54in. Simple
% one-term answers keep the tight default.
\newcommand{\eqline}[2][-0.5mm]{%
  \par\vspace{1mm}\noindent #2\par\vspace{#1}%
  \noindent\textcolor{rulehue}{\rule{\linewidth}{0.5pt}}\par\vspace{4mm}%
}

\begin{document}

\noteshead{Electric Potential Energy}{}

\goal{
\begin{itemize}[itemsep=5pt,parsep=2pt,topsep=2pt]
\item Build the idea of potential energy with gravity first: height above
the ground, and what happens if you keep going \emph{below} it.
\item See the same idea for a charge pushed against a uniform field ---
for both a $+$ and a $-$ charge.
\item Write the \emph{full} (non-uniform) form of gravitational PE, and see
that ``higher'' always means ``more PE,'' even though the value is negative.
\item Use that same point-source structure to build the electric PE of two
point charges, then a whole system of them.
\item Combine two point charges' PE to get the potential energy of a
dipole in a field.
\end{itemize}}

%============================================================
\nsec{Gravitational PE: height, then a hole}

A ball sits a height $h$ above the ground. Ground level is our reference:
$U=0$ there. Draw the ball, the ground, and then let it fall \emph{past}
ground level into a hole of depth $d$ --- draw that too.

\vspace{7cm}

\eqline{$U(h) = $}

Now let the ball fall past ground level, into a hole of depth $d$.

\eqline{$U(\text{bottom}) = $}

\eqline{$\Delta U = U(\text{bottom}) - U(h) = $}

\checkpoint{The ball's PE at the bottom of the hole came out negative. Is
that a problem?}

%============================================================
\nsec{The same idea, for a charge in a uniform field}

Replace gravity with a uniform electric field $\vec E$ pointing \emph{down}
(same visual sense as gravity), and try it with both signs of charge. Draw
a $+q$ and a $-q$ in the field, below, each with an arrow for the force
$q\vec E$ on it.

\vspace{6.5cm}

\begin{itemize}[itemsep=9pt,parsep=4pt,topsep=4pt]
\item Force on the $+q$: $q\vec E$ points \bl[6.5cm] --- that's its
\emph{natural} (free) direction of motion.
\item Force on the $-q$: $q\vec E$ points \bl[6.5cm].
\end{itemize}

Push each charge a distance $d$ \emph{against} its own natural direction ---
the same idea as raising the ball.

\eqline{$\Delta U_+ = $}
\eqline{$\Delta U_- = $}

\begin{itemize}[itemsep=9pt,parsep=4pt,topsep=4pt]
\item Both are \bl[2.4cm] --- pushing \emph{against} a charge's natural
direction always stores positive PE, regardless of the charge's sign.
\item Let either charge go instead, and it moves along its \emph{natural}
direction: PE \bl[2.4cm] either way --- the field does the work.
\end{itemize}

\checkpoint{A negative charge is released in this downward field. Which way
does it move, and what happens to its PE?}

%============================================================
\nsec{Gravitational PE, the full (non-uniform) form}

$U=mgh$ only works because $g$ barely changes over any height you can
reach near Earth's surface. Zoom out far enough to leave that ``uniform
field'' approximation, and the real, general form (valid at any separation
$r$ between two masses) is:

\eqline[1.4cm]{$U(r) = $}

Sketch $U$ vs.\ $r$ below.

\begin{center}
\begin{tikzpicture}[x=1cm,y=1cm]
  \draw[gray!70, ->] (0,-6.3) -- (0,0.6) node[above, scale=0.8] {$U$};
  \draw[gray!70, ->] (-0.3,0) -- (7.6,0) node[right, scale=0.8] {$r$};
  \draw[gray!50, dashed] (0,0) -- (7.3,0);
  \draw[gray!50] (1,0.08) -- (1,-0.08);
  \node[gray, scale=0.75] at (1,-0.4) {$R$};
  \node[gray, scale=0.7] at (-0.32,0.28) {$0$};
\end{tikzpicture}
\end{center}

\begin{itemize}[itemsep=9pt,parsep=4pt,topsep=4pt]
\item As $r\to\infty$, $U\to$ \bl[2.4cm] --- same rule as always: zero PE at
infinite separation.
\item Every value of $U(r)$ here is \bl[2.4cm]. But moving to a
\emph{larger} $r$ (further from Earth, ``higher'') always makes $U$
\bl[6.5cm]. Don't let the minus sign fool you: higher is \emph{always} more
PE, on this curve or any other.
\end{itemize}

\checkpoint{True or false: since $U(r)$ is always negative here, moving
\emph{further} from Earth makes $U$ smaller.}

%============================================================
\nsec{Electric PE of two point charges (from the class prep)}

Pulling a pair of point charges apart, or pushing them together, takes work
--- work that gets stored as potential energy, exactly like the point-mass
form above.

\eqline[1.4cm]{$U(q_1,q_2) = $}

Sketch $U$ vs.\ $r_{12}$ below, for an opposite-sign pair and a like-sign
pair, on the same axes. (Same shape you just drew for gravity --- plus its
mirror image.)

\begin{center}
\begin{tikzpicture}[x=1cm,y=1cm]
  \draw[gray!70, ->] (0,-3.6) -- (0,3.6) node[above, scale=0.8] {$U$};
  \draw[gray!70, ->] (-0.3,0) -- (8.6,0) node[right, scale=0.8] {$r_{12}$};
  \draw[gray!50, dashed] (0.55,0) -- (7.3,0);
\end{tikzpicture}
\end{center}

\begin{itemize}[itemsep=9pt,parsep=4pt,topsep=4pt]
\item As $r_{12}\to\infty$, $U\to$ \bl[2.4cm].
\item For three or more charges, sum over every \bl[3.2cm] (not every charge).
\item Opposite charges ($+$ and $-$): $U$ is \bl[2.4cm]\ at any finite
separation, and $U\to0$ only as $r\to\infty$. So $U=0$ is actually the
\bl[3.2cm]\ value here --- pulling the charges closer always finds a
\emph{lower} (more negative) PE.
\item Like charges (same sign): $U$ is \bl[2.4cm]\ at any finite separation,
and again $U\to0$ as $r\to\infty$. Here $U=0$ \emph{is} the \bl[3.2cm]\
value --- you can't push like charges to a lower PE than infinitely far
apart.
\end{itemize}

\eqline[1.7cm]{$U_{\text{total}} = $}

\checkpoint{Two $+q$ charges are brought from far apart to a separation $r$.
Is $U$ positive or negative? What about a $+q$ and a $-q$?}

%============================================================
\nsec{The dipole in a field}

A dipole is a $+q$ and a $-q$ separated by $\vec d$ (dipole moment
$\vec p = q\vec d$), sitting in a uniform field $\vec E$, tilted at angle
$\theta$ between $\vec p$ and $\vec E$. Each charge sits at a different
place along $\vec E$ --- use \S2's pushing idea on each charge, then add.

\vspace{7cm}

\eqline{$U(\theta) = $}

\begin{itemize}[itemsep=9pt,parsep=4pt,topsep=4pt]
\item At $\theta=90^\circ$ (perpendicular to the field): $U=$ \bl[2.4cm].
\item At $\theta=0^\circ$ (aligned with the field): $U=$ \bl[3.2cm] ---
which is \bl[2.4cm]\ than the $\theta=90^\circ$ value!
\item At $\theta=180^\circ$ (anti-aligned): $U=$ \bl[3.2cm].
\end{itemize}

\checkpoint{Which orientation is the most stable (lowest PE)? Is that the
orientation where $U=0$? What does this have in common with the ball in
the hole?}

%============================================================
\nsec{Sketch it: $U$ vs.\ $\theta$}

\begin{center}
\begin{tikzpicture}[x=1cm,y=1cm]
  \draw[gray!70, ->] (0,-2.1) -- (0,2.3) node[above, scale=0.8] {$U$};
  \draw[gray!70, ->] (0,0) -- (9.2,0) node[right, scale=0.8] {$\theta$};
  \foreach \x/\lab in {0/{$0^\circ$}, 2.2/{$90^\circ$}, 4.4/{$180^\circ$}, 6.6/{$270^\circ$}, 8.8/{$360^\circ$}}{
    \draw[gray!50] (\x,0.08) -- (\x,-0.08);
    \node[gray, scale=0.75] at (\x,-0.4) {\lab};
  }
\end{tikzpicture}
\end{center}

{\small\color{blankink} Check your curves, and the point-charge and
uniform-field results above, against the \emph{Assembling Charges},
\emph{PE vs.\ Separation}, and \emph{Electric Dipole in a Field} sims ---
open their \textbf{Info} panels.}

\end{document}
